\section{10. Running traps}

The dungeon contains door traps, ordinary room traps and death traps. Their purposes and severity should differ.

\subsection{Door traps}

Door traps punish careless interaction with entrances and exits. They should usually cause limited damage, a temporary condition or an immediate complication rather than kill a healthy character outright. Some trigger on opening, others when a handle is touched or a threshold crossed.

For example, a Needle Lock might release a poisoned needle when touched, while a Flame Seal might release a burst of fire toward the opener. Use the trap description to establish the trigger and saving throw. Once triggered or disarmed, the active trap is removed from the corresponding door state. Its appearance remains recognizable.

\subsection{Ordinary room traps}

Ordinary traps are tricky rooms containing a reward chest. They present obstacles characters can navigate, disable or endure: dart launchers, pendulums, electrified floors, steam vents, grasping vegetation and slippery surfaces. Solutions might involve checks, saving throws, equipment, spells or clever movement. Characters need not permanently disable the hazard before retrieving its chest; a daring character may reach the reward and escape.

\subsection{Death traps}

Death traps are dangerous rooms without ordinary reward chests. They create urgent survival situations: crushing walls, rising water, collapsing floors, incinerators, petrification and disintegration. They should be capable of killing characters who remain exposed or repeatedly fail to escape, but should not normally kill a healthy character without warning or opportunity to react.

\subsection{A useful escalation pattern}

1. Warning: the mechanism activates and presents an unmistakable threat.

2. Escalation: the hazard intensifies; movement, checks and countermeasures matter.

3. Critical danger: severe consequences for those who remain.

\subsection{Suggested damage scaling}

\begin{longtable}{p{0.34\linewidth}p{0.6\linewidth}}

\rowcolor{tablehead} \textbf{Hazard} & \textbf{Suggested starting damage (levels 5–8)} \\

\endfirsthead

\rowcolor{tablehead} \textbf{Hazard} & \textbf{Suggested starting damage (levels 5–8)} \\

\endhead

Minor door trap & 2d6–3d6 \\

Ordinary room trap & 3d6–5d6 \\

Severe room hazard & 5d6–8d6 \\

Critical death trap & 8d6 or more, as appropriate \\

\end{longtable}

Adjust for frequency, area, countermeasures and whether successful saves reduce or prevent damage. A repeating hazard damaging every character each round can be far more dangerous than a single high-damage effect.
